% ---------------------------------------------------------------------------
% Appendix C.5 — Persona Format Ablation (Structured vs. Natural-Language)
% Appendix C'nin (Calibration Funnel: Cell-Level Results) SONUNA, C.4'ün
% ardından beşinci alt bölüm olarak yapıştırılır.
% Veri: outputs/natural_ablation/metrics_side_by_side.csv
% Figür: outputs/natural_ablation/fig_structured_vs_natural.pdf
%
% ANA METİN BAĞLANTISI — §3.4'teki revize "Prompt-framing stress tests"
% paragrafı bu ablasyonu dördüncü varyant olarak tanıtır ve
% \ref{app:natural}'a footnote ile işaret eder; §4.1'deki stres testi bulgu
% paragrafına da tek cümlelik özet eklenir (calibration_revision_notes'a
% bakınız).
%
% NOT: pacdemons hücresi T=0.8 ile koşulmuştur; nihai protokol T=0 olarak
% raporlanacaksa ablasyonun T=0.8'de yapıldığı dipnotta belirtilmeli ya da
% natural koşusu T=0 ile tekrarlanmalıdır.
% ---------------------------------------------------------------------------

\subsection{Persona Format Ablation: Structured versus Natural-Language
Profiles}
\label{app:natural}

The persona profiles used throughout the paper are structured: each survey
variable appears as a labeled line under a fixed group heading
(Appendix~D). A plausible alternative is to present the same information as
flowing prose. This appendix tests whether that choice matters.

\paragraph{Procedure.} For every respondent in the selected C6
configuration, GPT-4o-mini rewrites the structured profile as a third-person
Turkish biography under strict instructions: preserve every recorded fact,
add no new information, add no interpretation or evaluative voice, and omit
rather than paraphrase missing values. The rewritten biographies replace the
structured profiles in the two final calibrated cells, with every other
component of the inference setup held fixed; the structured rows below are
therefore the final-protocol results from the main text.

\begin{table}[h]
\centering
\caption{Structured versus natural-language personas under the final
calibrated protocols. Arrows indicate the direction of better performance.
Rec(+) denotes protest-positive recall; $\kappa_w$ is quadratic weighted
kappa.}
\label{tab:natural}
\begin{tabular}{llccccc}
\hline
Outcome & Persona format & JSD $\downarrow$ & Wass.\ $\downarrow$ &
MCC $\uparrow$ & Rec(+) $\uparrow$ & $\kappa_w$ $\uparrow$ \\
\hline
\emph{pacdemons} & Structured & 0.001 & 0.017 & 0.430 & 0.536 & --- \\
                 & Natural    & 0.004 & 0.037 & 0.257 & 0.405 & --- \\[2pt]
\emph{womenwork} & Structured & 0.030 & 0.218 & ---   & ---   & 0.401 \\
                 & Natural    & 0.020 & 0.260 & ---   & ---   & 0.165 \\
\hline
\end{tabular}
\end{table}

\paragraph{Results.} Aggregate distributional fit is broadly preserved and
not uniformly ordered: the natural format slightly worsens both
distributional metrics for \emph{pacdemons} and moves them in opposite
directions for \emph{womenwork}, improving Jensen--Shannon divergence while
worsening Wasserstein distance. Respondent-level agreement, by contrast,
degrades sharply and consistently. On the rare binary outcome, the Matthews
correlation coefficient falls from 0.43 to 0.26 and protest-positive recall
from 53.6 to 40.5 percent. On the ordered attitudinal outcome, quadratic
weighted kappa falls from 0.40 to 0.17, less than half its structured
value. Figure~\ref{fig:natural} displays the full comparison.

\begin{figure}[h]
  \centering
  \includegraphics[width=\textwidth]{figures/fig_structured_vs_natural.pdf}
  \caption{Structured versus natural-language personas across all
  calibration metrics, for both outcomes, under the final calibrated
  protocols.}
  \label{fig:natural}
\end{figure}

\paragraph{Interpretation.} The structured format gives the model an
explicit map of where each attribute begins and ends and what value it
takes; the biographical rewrite embeds the same facts in sentences and
renders scaled quantities as partial verbal cues. What is lost in that
translation is evidently not the shape of the aggregate distribution, which
either format can approximate, but the respondent-conditional signal the
model needs to infer how a specific individual answers. The ablation
therefore supports the design decision of Section~3.3 on empirical rather
than only conceptual grounds, and it illustrates why calibration is
evaluated at two levels throughout the paper: a persona representation can
pass a purely distributional audit while failing the respondent-level one.
